% GENERATED FILE. Edit canonical lessons, then run npm run generate:books.
% canonical-source-hash: fnv1a64:1e13670b39f372ed
% canonical-lessons: BN-C500-bibhag, BN-C500-hukum, BN-C500-kormosthol, BN-C500-sara, BN-C500-todarok-kora

\chapter{Orders and Departments}
\label{ch:bn-orders-and-departments}
\hlchaptermodality{500}

\begin{chapteropening}
\textbf{By the end of this chapter:} \emph{I can talk about departments, instructions and getting work done.}

Complete the last lesson of chapter 500: I can talk about departments, instructions and getting work done.
\end{chapteropening}

\section[\texorpdfstring{\bn{বিভাগ} (bibhāg)}{বিভাগ (bibhāg)}]{\bn{বিভাগ} (bibhāg) — a department; a division}

\label{lesson:BN-C500-bibhag}



\begin{quote}
Before the new one: say the Bengali for a supply, then the Bengali for imports.
\end{quote}

\subsection*{You'll want to know: \bn{বিভাগ}}
\textbf{\bn{বিভাগ}} — \emph{bibhāg} — \textquotedblleft{}a department; a division\textquotedblright{}.

\begin{quote}
\bn{আমার} \bn{বোন} \bn{এই} \bn{কোম্পানির} \bn{নতুন} \bn{বিভাগে} \bn{কাজ} \bn{করে।}
\end{quote}

\begin{quote}
\emph{Āmār bon ei kompānir nôtun bibhāge kāj kôre.}
\end{quote}

\begin{quote}
My sister works in the company's new department.
\end{quote}

\begin{cousinweb}
\textbf{\bn{বিভাগ}} (\emph{bibhāg}) is a tatsama, Sanskrit \emph{vibhāga}, \textquotedblleft{}distribution, a share\textquotedblright{} (Monier-Williams): a part of the whole.
\end{cousinweb}

\begin{grammarlens}[title={What you've built}]
One more word for describing what you lived and felt.
\end{grammarlens}

\subsection*{Your turn}
\begin{itemize}
\raggedright
  \item \emph{Say it:} \emph{bibhāg}
  \item \emph{Say it:} \emph{bibhāg}, once more
  \item \emph{Say it:} \emph{Āmār bon ei kompānir nôtun bibhāge kāj kôre.} — My sister works in the company's new department.
  \item \emph{Recall:} say the Bengali for a supply, then the Bengali for imports, then say \emph{bibhāg} again
\end{itemize}

\begin{tcolorbox}[breakable,colback=teal!4,colframe=teal!35!black,title={Before you move on}]
What does \bn{বিভাগ} mean? (\textquotedblleft{}A department; a division\textquotedblright{}.) Say it once more.
\end{tcolorbox}
\section[\texorpdfstring{\bn{হুকুম} (hukum)}{হুকুম (hukum)}]{\bn{হুকুম} (hukum) — an order; a command}

\label{lesson:BN-C500-hukum}



\begin{quote}
Before the new one: say the Bengali for imports, then the Bengali for a department.
\end{quote}

\subsection*{You'll want to know: \bn{হুকুম}}
\textbf{\bn{হুকুম}} — \emph{hukum} — \textquotedblleft{}an order; a command\textquotedblright{}.

\begin{quote}
\bn{ম্যানেজার} \bn{নতুন} \bn{হুকুম} \bn{দিলেন।}
\end{quote}

\begin{quote}
\emph{Mænejār nôtun hukum dilen.}
\end{quote}

\begin{quote}
The manager gave a new order.
\end{quote}

\begin{cousinweb}
\textbf{\bn{হুকুম}} (\emph{hukum}) is usually traced to Arabic \emph{ḥukm}, \textquotedblleft{}judgement, a command\textquotedblright{}, through Persian.
\end{cousinweb}

\begin{grammarlens}[title={What you've built}]
One more word for describing what you lived and felt.
\end{grammarlens}

\subsection*{Your turn}
\begin{itemize}
\raggedright
  \item \emph{Say it:} \emph{hukum}
  \item \emph{Say it:} \emph{hukum}, once more
  \item \emph{Say it:} \emph{Mænejār nôtun hukum dilen.} — The manager gave a new order.
  \item \emph{Recall:} say the Bengali for imports, then the Bengali for a department, then say \emph{hukum} again
\end{itemize}

\begin{tcolorbox}[breakable,colback=teal!4,colframe=teal!35!black,title={Before you move on}]
What does \bn{হুকুম} mean? (\textquotedblleft{}An order; a command\textquotedblright{}.) Say it once more.
\end{tcolorbox}
\section[\texorpdfstring{\bn{কর্মস্থল} (kôrmôsthôl)}{কর্মস্থল (kôrmôsthôl)}]{\bn{কর্মস্থল} (kôrmôsthôl) — a workplace}

\label{lesson:BN-C500-kormosthol}



\begin{quote}
Before the new one: say the Bengali for a department, then the Bengali for a command.
\end{quote}

\subsection*{You'll want to know: \bn{কর্মস্থল}}
\textbf{\bn{কর্মস্থল}} — \emph{kôrmôsthôl} — \textquotedblleft{}a workplace\textquotedblright{}.

\begin{quote}
\bn{আমার} \bn{কর্মস্থল} \bn{বাড়ি} \bn{থেকে} \bn{দূরে।}
\end{quote}

\begin{quote}
\emph{Āmār kôrmôsthôl bāṛi theke dūre.}
\end{quote}

\begin{quote}
My workplace is far from home.
\end{quote}

\begin{cousinweb}
\textbf{\bn{কর্মস্থল}} (\emph{kôrmôsthôl}) joins Sanskrit \emph{karman}, \textquotedblleft{}work\textquotedblright{}, and \emph{sthala}, \textquotedblleft{}a place, ground\textquotedblright{} (Monier-Williams).
\end{cousinweb}

\begin{grammarlens}[title={What you've built}]
One more word for describing what you lived and felt.
\end{grammarlens}

\subsection*{Your turn}
\begin{itemize}
\raggedright
  \item \emph{Say it:} \emph{kôrmôsthôl}
  \item \emph{Say it:} \emph{kôrmôsthôl}, once more
  \item \emph{Say it:} \emph{Āmār kôrmôsthôl bāṛi theke dūre.} — My workplace is far from home.
  \item \emph{Recall:} say the Bengali for a department, then the Bengali for a command, then say \emph{kôrmôsthôl} again
\end{itemize}

\begin{tcolorbox}[breakable,colback=teal!4,colframe=teal!35!black,title={Before you move on}]
What does \bn{কর্মস্থল} mean? (\textquotedblleft{}A workplace\textquotedblright{}.) Say it once more.
\end{tcolorbox}
\section[\texorpdfstring{\bn{সারা} (sārā)}{সারা (sārā)}]{\bn{সারা} (sārā) — to finish; to get done}

\label{lesson:BN-C500-sara}



\begin{quote}
Before the new one: say the Bengali for a command, then the Bengali for a workplace.
\end{quote}

\subsection*{You'll want to know: \bn{সারা}}
\textbf{\bn{সারা}} — \emph{sārā} — \textquotedblleft{}to finish; to get done\textquotedblright{}.

\begin{quote}
\bn{আমি} \bn{দুপুরের} \bn{আগে} \bn{কাজটা} \bn{সারলাম।}
\end{quote}

\begin{quote}
\emph{Āmi dupurer āge kājṭā sārlām.}
\end{quote}

\begin{quote}
I got the work done before midday.
\end{quote}

Past (he/she …): \bn{তিনি} \bn{রান্না} \bn{সারলেন}, \emph{tini rānnā sārlen}.

\begin{cousinweb}
\textbf{\bn{সারা}} (\emph{sārā}) is a tadbhava from Sanskrit \emph{sārayati}, \textquotedblleft{}makes move, removes\textquotedblright{}; CDIAL gives Bengali \emph{sārā} both for repairing and for finishing (CDIAL 13358).
\end{cousinweb}

\begin{grammarlens}[title={What you've built}]
One more word for describing what you lived and felt.
\end{grammarlens}

\subsection*{Your turn}
\begin{itemize}
\raggedright
  \item \emph{Say it:} \emph{sārā}
  \item \emph{Say it:} \emph{sārā}, once more
  \item \emph{Say it:} \emph{Āmi dupurer āge kājṭā sārlām.} — I got the work done before midday.
  \item \emph{Recall:} say the Bengali for a command, then the Bengali for a workplace, then say \emph{sārā} again
\end{itemize}

\begin{tcolorbox}[breakable,colback=teal!4,colframe=teal!35!black,title={Before you move on}]
What does \bn{সারা} mean? (\textquotedblleft{}To finish; to get done\textquotedblright{}.) Say it once more.
\end{tcolorbox}
\section[\texorpdfstring{\bn{তদারক} \bn{করা} (tôdārôk kôrā)}{তদারক করা (tôdārôk kôrā)}]{\bn{তদারক} \bn{করা} (tôdārôk kôrā) — to supervise}

\label{lesson:BN-C500-todarok-kora}



\begin{quote}
Before the new one: say the Bengali for a workplace, then the Bengali for to get done.
\end{quote}

\subsection*{You'll want to know: \bn{তদারক} \bn{করা}}
\textbf{\bn{তদারক} \bn{করা}} — \emph{tôdārôk kôrā} — \textquotedblleft{}to supervise\textquotedblright{}.

\begin{quote}
\bn{নতুন} \bn{ম্যানেজার} \bn{কারখানার} \bn{কাজ} \bn{তদারক} \bn{করেন।}
\end{quote}

\begin{quote}
\emph{Nôtun mænejār kārkhānār kāj tôdārôk kôren.}
\end{quote}

\begin{quote}
The new manager supervises the work at the factory.
\end{quote}

Past (he/she …): \bn{তিনি} \bn{কাজটা} \bn{তদারক} \bn{করলেন}, \emph{tini kājṭā tôdārôk kôrlen}.

\begin{cousinweb}
\textbf{\bn{তদারক} \bn{করা}} (\emph{tôdārôk kôrā}) joins \bn{তদারক}, \textquotedblleft{}supervision\textquotedblright{}, to \bn{করা}; \bn{তদারক} is usually traced to Arabic \emph{tadāruk}, \textquotedblleft{}setting right\textquotedblright{}, through Persian.
\end{cousinweb}

\begin{grammarlens}[title={What you've built}]
One more word for describing what you lived and felt.
\end{grammarlens}

\subsection*{Your turn}
\begin{itemize}
\raggedright
  \item \emph{Say it:} \emph{tôdārôk kôrā}
  \item \emph{Say it:} \emph{tôdārôk kôrā}, once more
  \item \emph{Say it:} \emph{Nôtun mænejār kārkhānār kāj tôdārôk kôren.} — The new manager supervises the work at the factory.
  \item \emph{Recall:} say the Bengali for a workplace, then the Bengali for to get done, then say \emph{tôdārôk kôrā} again
\end{itemize}

\begin{tcolorbox}[breakable,colback=teal!4,colframe=teal!35!black,title={Before you move on}]
What does \bn{তদারক} \bn{করা} mean? (\textquotedblleft{}To supervise\textquotedblright{}.) Say it once more.
\end{tcolorbox}
